
\documentclass[12pt]{article} % use larger type; default would be 10pt
\usepackage[square]{natbib}
\usepackage[utf8]{inputenc} % set input encoding (not needed with XeLaTeX)
\usepackage{amsmath,amssymb}
\DeclareMathOperator{\E}{\mathbb{E}}
\DeclareMathOperator{\sigw}{\sigma_w^2}

\usepackage{color}
\usepackage{pdfpages}
%%% PAGE DIMENSIONS
\usepackage{geometry} % to change the page dimensions
%\geometry{a4paper} % or letterpaper (US) or a5paper or....
\geometry{margin=1in} % for example, change the margins to 2 inches all round
% \geometry{landscape} % set up the page for landscape
%   read geometry.pdf for detailed page layout information
 \newcommand{\eps}{\varepsilon}
\usepackage{graphicx} % support the \includegraphics command and options

% \usepackage[parfill]{parskip} % Activate to begin paragraphs with an empty line rather than an indent

%%% PACKAGES
\usepackage{booktabs} % for much better looking tables
\usepackage{array} % for better arrays (eg matrices) in maths
\usepackage{paralist} % very flexible & customisable lists (eg. enumerate/itemize, etc.)
\usepackage{verbatim} % adds environment for commenting out blocks of text & for better verbatim
\usepackage{subfig} % make it possible to include more than one captioned figure/table in a single float
% These packages are all incorporated in the memoir class to one degree or another...

%%% HEADERS & FOOTERS
\usepackage{fancyhdr} % This should be set AFTER setting up the page geometry
\pagestyle{fancy} % options: empty , plain , fancy
\fancyhf{}
\renewcommand{\headrulewidth}{1pt} % customise the layout...
\lhead{STAT 429 Time Series Analysis}\chead{FALL 2020}\rhead{Instructor: Hyoeun Lee}
\lfoot{}\cfoot{}\rfoot{Page \thepage}




%%% SECTION TITLE APPEARANCE
\usepackage{sectsty}
\allsectionsfont{\sffamily\mdseries\upshape} % (See the fntguide.pdf for font help)
% (This matches ConTeXt defaults)

%%% ToC (table of contents) APPEARANCE
\usepackage[nottoc,notlof,notlot]{tocbibind} % Put the bibliography in the ToC
\usepackage[titles,subfigure]{tocloft} % Alter the style of the Table of Contents
\renewcommand{\cftsecfont}{\rmfamily\mdseries\upshape}
\renewcommand{\cftsecpagefont}{\rmfamily\mdseries\upshape} % No bold!


\definecolor{Red}{rgb}{1,0,0}
\newcommand{\Red}{\color{Red}}
\definecolor{Blue}{rgb}{0,0,1}
\newcommand{\Blue}{\color{Blue}}


%%%%%%%THEOREM 
\usepackage{thmtools,thm-restate}
\newtheorem{theorem}{Theorem}
\newtheorem{lemma}[theorem]{Lemma}
\newtheorem{proposition}[theorem]{Proposition}
\newtheorem{corollary}[theorem]{Corollary}
\newtheorem{definition}[theorem]{Definition}

\newenvironment{proof}[1][Proof]{\begin{trivlist}
\item[\hskip \labelsep {\bfseries #1}]}{\end{trivlist}}
%\newenvironment{definition}[1][Definition]{\begin{trivlist}
%\item[\hskip \labelsep {\bfseries #1}]}{\end{trivlist}}
\newenvironment{example}[1][Example]{\begin{trivlist}
\item[\hskip \labelsep {\bfseries #1}]}{\end{trivlist}}
\newenvironment{remark}[1][Remark]{\begin{trivlist}
\item[\hskip \labelsep {\bfseries #1}]}{\end{trivlist}}

\newcommand{\qed}{\nobreak \ifvmode \relax \else
      \ifdim\lastskip<1.5em \hskip-\lastskip
      \hskip1.5em plus0em minus0.5em \fi \nobreak
      \vrule height0.75em width0.5em depth0.25em\fi}
   
   
   
%   \usepackage{nopageno}   
 %%%%%%%%%%%%%%

%\linespread{2}





%%% END Article customizations



%%% The "real" document content comes below...


\date{}


%\title{Chapter 1}

\begin{document}




\textbf{\Large 2.1  Correlation and Stationary Time Series: \\ Measuring Dependence}
%\textbf{\Large Chapter 1 Time Series Elements}\\
%\textbf{\large 1.1 Introduction}

\medskip


%\noindent
\fbox{%
    \parbox{0.9\textwidth}{%
 


\textbf{Learning outcome of Section 2.1}

\begin{itemize}
\item Understand mean function of time series and its relationship to the trend
\item Understand autocovariance, autocorrelation function and its relationship to volatility and dependence
\item Understand cross-covariance, cross-correlation function and its role
\end{itemize}



    }%
}




\medskip
\medskip

\noindent\textbf{Recall}

Time Series: 
\begin{itemize}
	\item a Discrete Stochastic Process
	\item Denote: $\{X_t\}_{t\geq 0}$, $t$: time index
	\item ${t\geq 0}$ or $ t=0, \pm 1, \pm 2, \pm 3, \dots $ (depends on starting point)
	\item 3 things important for modeling:  Mean, Variance, Dependence
\end{itemize}

\noindent \textbf{Properties of the Mean, Variance and Covariance}
\begin{itemize}
\item $\E[aX+bY+c]=a\E[X]+b\E[Y]+c$
\item $Var(X)=\E[(X-\E[X])^2]=\E[X^2]-(\E[X])^2$
\item $Var(a+bX)=b^2Var(X)$
\item If $X$ and $Y$ are independent $Var(X+Y)=Var(X)+Var(Y)$
\item $Cov(X,Y)=\E[(X-\E[X])(Y-\E[Y])]$ 
\item  $Cov(a+bX,c+dY)=bd Cov(X,Y)$
\item $Var(X+Y)=Var(X)+Var(Y)+2Cov(X,Y)$
\item $Cov(X+Y,Z)=Cov(X,Z)+Cov(Y,Z)$
\item $Cov(X,X)=Var(X)$
\item $Cov(X,Y)=Cov(Y,X)$
\end{itemize}


\newpage

\noindent\textbf{Mean Function}
$$\mu_X(t)= \mathbb{E}[X_t]$$


\begin{enumerate}
	\item $X_t = \frac{1}{2} (W_{t-1} + W_t )$
	\item[] $\mu_{X} (t)=0$ 
	\item $X_t=\delta t + \sum_{j=1}^{t} W_j$
	\item[] $\mu_{X} (t)= \delta t$
	\item $X_t = 2 \cos (2\pi \frac{t+15}{50}) + W_t$
	\item[] $\mu_{X} (t)= 2 \cos (2\pi \frac{t+15}{50})$
\end{enumerate}

\medskip


\noindent\textbf
{Autocovariance function}

\begin{align*}
	\gamma_X ({t,s})=Cov(X_t,X_s)=&\E[(X_t-\mu_X(t))(X_s-\mu_X(s))]\\
	=&\E[X_t X_s]-\mu_X(t) \mu_X(s)
\end{align*}





\begin{enumerate}
	\item Autocovariance of White Noise:
	\begin{equation}
		\gamma_{W} (t,s)= Cov (W_t, W_s) = 
		\begin{cases}
			\sigma_W^2 & s=t\\
			0 & s\neq t.
		\end{cases}
	\end{equation}

	\item $X_t = \frac{1}{2} (W_{t-1} + W_t )$
	\item[] 
	\begin{equation}
		\gamma_{X} (t,s)=\begin{cases}
			\frac{1}{2} \sigw & s=t,\\
			\frac{1}{4}	\sigw & |t-s|=1,\\
			0 & |t-s|>1.
		\end{cases}
	\end{equation}
	\item $X_t=\delta t + \sum_{j=1}^{t} W_j$
	\item[] 
	\begin{equation}
		\gamma_{X} (t,s)= Cov \left( \delta t + \sum_{j=1}^{t} W_j, \delta s + \sum_{j=1}^{s} W_j\right)
		=\min(s,t) \sigw
	\end{equation}
	\item $X_t = 2 \cos (2\pi \frac{t+15}{50}) + W_t$
	\item[] 
	\begin{equation}
		\gamma_{X} (t,s)=\sigw
	\end{equation}
\end{enumerate}

\medskip


\noindent\textbf
{Autocorrelation function}

$$\rho_{t,s}= Corr(Y_t,Y_s)$$
$ t,s=0, \pm 1, \pm 2, \pm 3, \dots $
where $$Corr(Y_t,Y_s)=\dfrac{Cov(Y_t,Y_s)}{\sqrt{Var(Y_t) Var(Y_s)}}=
\dfrac{\gamma_{t,s}}{\sqrt{\gamma_{t,t }\gamma_{s,s}}} $$

Covariance and correlation are a measure of \bf{linear} relationship between two random variables.


\newpage
{Important properties of the autocorrelation and autocovariance function}
\begin{table}
\begin{tabular}{cc}
$\gamma_{t,t}=Var(Y_t)$ & $\rho_{t,t}=1$\\ \\
$\gamma_{t,s}=\gamma_{s,t}$ & $\rho_{t,s}=\rho_{s,t}$\\ \\
$|\gamma_{t,s}|\leq \sqrt{\gamma_{t,t} \gamma_{s,s}}$ & $|\rho_{t,s}|\leq 1$
\end{tabular}
\end{table}

%
%{Assignments}
%Attempt the following problems:
%\begin{enumerate} 
%\item Use the Cauchy $–$ Schwarz inequality to test the property:
%$$|\gamma_{t,s}|\leq \sqrt{\gamma_{t,t }\gamma_{s,s}}$$
%\item Proof that: 
%$$ Cov[\Sigma_{i=1}^m c_i Y_{t_i},\Sigma_{j=1}^n d_j Y_{s_j} ]=
%\Sigma_{i=1}^m \Sigma_{j=1}^n c_i d_j Cov(Y_{t_i},Y_{s_j}) $$
%for constants $c_1,c_2,\ldots, c_m$, $d_1,d_2,\ldots, d_n$ and time points $t_1,t_2,\ldots,t_m$ and $s_1,s_2,\ldots,s_n$
%\item Obtain an expression for $$Var[\Sigma_{i=1}^n c_iY_{t_i}]$$ as a special case
%\item Assume: $E(X)=4$,$Var(X)=11$,$E(Y)=0$,$Var(Y)=5$, and $Corr(X,Y)=0.3$. Find:
%\begin{itemize}
%\item $Var(X+Y)$
%\item $Cov(X,X+Y)$
%\item $Corr(X+Y,X-Y)$
%\end{itemize}
%\end{enumerate}

%
%{Assignments (Cont.)}
%\only<1>{
%Use \textsf{R} for the following problems: 
%\begin{enumerate} 
%\item Plot the Time Series EQ5 and EXP6 of package {\bf astsa} in a similar format as example 6 in lecture notes Week 1 Class 1 (use \text{help(EQ5)} and \text{help(EXP6)} for more information on the data sets)
%\item Plot the two data sets on the same graph using different colors or different line types and comment on the results
%\item Consider a signal-plus-noise model of the general form $x_t = s_t + w_t$,
%where $w_t$ is a Gaussian white noise with $\sigma^2_w=1$. Simulate and plot $n = 200$
%observations from each of the following two models:
%\end{enumerate}}
%\only<2>{
%\begin{enumerate}[a]
%\item $x_t= s_t + w_t$, for $t=1,2,\ldots,200$, where:
%\[
%s_t= \left\{
%\begin{array}{ll}
%0, & t=1,\ldots, 100\\
%10\exp\{-\frac{(t-100)}{20}\}\cos(2\pi t/4), & t=101,\ldots, 200
%\end{array}
%\right.
%\]
%\item $x_t= s_t + w_t$, for $t=1,2,\ldots,200$, where:
%\[
%s_t= \left\{
%\begin{array}{ll}
%0, & t=1,\ldots, 100\\
%10\exp\{-\frac{(t-100)}{200}\}\cos(2\pi t/4), & t=101,\ldots, 200
%\end{array}
%\right.
%\]
%\item Compare the general appearance of series (a) and (b) with the Time Series plotted in 1). In addition plot or sketch the signal modulators $\exp\{-t/20\}$ and  $\exp\{-t/200\}$
%\end{enumerate}
%}
%
%
%[fragile]
%Hint for the last problem:
%\begin{verbatim}
%s = c(rep(0,100), 10*exp(-(1:100)/20)*cos(2*pi*1:100/4))
%x = ts(s + rnorm(200, 0, 1))
%\end{verbatim}







\hrulefill

\medskip

Next: 


\hrulefill

Next: 2.2 



\end{document}
